\section{原群代数,幺半群代数,群代数}

记号约定:$A,B$是交换环,$S,S^{\prime}$是原群(乘法记号),$E=A^{\left(S\right)}$.

\begin{definition}{原群代数}{}
	设$\left(e_{s}\right)_{s\in S}$是$A^{\left(S\right)}$的典范基,定义$A^{\left(S\right)}$上的乘法$\cdot$使$\forall s,t\in S\left(e_{s}\cdot e_{t}=e_{st}\right)$,则$A^{\left(S\right)}$是$A$-代数,称为$S$在$A$上的原群代数.
\end{definition}

\begin{remark}{}{}
	设$x,y\in A^{\left(S\right)}$,其中$x=\sum\limits_{s\in S}\alpha_{s}e_{s},y=\sum\limits_{s\in S}\beta_{s}e_{t}$,则$xy=\sum\limits_{s\in S}\left(\sum\limits_{\substack{t,u\in S\\tu=s}}\alpha_{t}\beta_{u}\right)e_{s}$.
\end{remark}

\begin{definition}{幺半群代数,群代数}{}
	若$S$是幺半群(相对的,群),则称$S$在$A$上的原群代数$A^{\left(S\right)}$为幺半群(相对的,群)代数.
\end{definition}

\begin{proposition}{幺半群代数和群代数的性质}{}
	\begin{enumerate}
		\item 若$S$是幺半群(相对的,群),则$A^{\left(S\right)}$是结合代数.
		\item 若$S$是交换幺半群,则$A^{\left(S\right)}$是交换结合代数.
		\item 若$u$是$S$的幺元,则$e_{u}$是$A^{\left(S\right)}$的代数,并且$A$典范同构于$Ae_{u}$.
	\end{enumerate}
\end{proposition}

\begin{remark}{记号说明}{}
	当$A$是非零环时,常把$S$和典范映射$S\to A^{\left(S\right)},s\mapsto e_{s}$的像等同起来.由此,$A^{\left(S\right)}$的每个元素\textbf{典范等同}于$\sum\limits_{s\in S}\alpha_{s}s$.需要指出的是,当$S$采用加法记号时,这种等同会造成混淆.为了消除歧义,我们会把$e_{s}$写作$e^{s}$.
\end{remark}

\begin{proposition}{标量扩张下的同构}{}
	设$\rho\in\Hom_{\mathsf{Ring}}\left(A,B\right)$,则存在唯一的典范同构$j\in\Hom_{B-\mathsf{Alg}}\left(\rho^{\ast}\left(A^{\left(S\right)}\right),B^{\left(S\right)}\right)$使
	\begin{align*}
		\forall s\in S\left(j\left(\rho^{\ast}\left(e_{s}\right)\right)=e_{s}^{\prime}\right),
	\end{align*}
	其中$\left(e_{s}\right)_{s\in S},\left(e_{s}^{\prime}\right)_{s\in S}$分别是$A^{\left(S\right)},B^{\left(S\right)}$的典范基.
\end{proposition}

\begin{proposition}{原群代数的泛性质}{}
	设$F$是$A$-代数.对任一$f\in\Hom_{\mathsf{Mag}}\left(S,F\right)$,存在唯一的$\bar{f}\in\Hom_{A-\mathsf{Alg}}\left(A^{\left(S\right)},F\right)$使下图交换.
	\begin{center}
		\begin{tikzcd}
			S && F \\
			\\
			{A^{\left(S\right)}} && F
			\arrow["f", from=1-1, to=1-3]
			\arrow[ hook, from=1-1, to=3-1]
			\arrow["{\id_{F}}", from=1-3, to=3-3]
			\arrow["{\exists!\bar{f}}"', dashed, from=3-1, to=3-3]
		\end{tikzcd}
	\end{center}
	其中,水平方向的映射$S\to A^{\left(S\right)}$是典范嵌入.
\end{proposition}

\begin{corollary}{原群代数的函子性}{}
	设$f\in\Hom_{\mathsf{Mag}}\left(S,S^{\prime}\right)$,则存在唯一的$\bar{f}\in\Hom_{A-\mathsf{Alg}}\left(A^{\left(S\right)},A^{\left(S^{\prime}\right)}\right)$使下图交换.
	\begin{center}
		\begin{tikzcd}
			S && {S^{\prime}} \\
			\\
			{A^{\left(S\right)}} && {A^{S^{\prime}}}
			\arrow["f", from=1-1, to=1-3]
			\arrow[hook, from=1-1, to=3-1]
			\arrow[hook', from=1-3, to=3-3]
			\arrow["{\exists!\bar{f}}"', dashed, from=3-1, to=3-3]
		\end{tikzcd}
	\end{center}
	其中,水平方向的映射$S\to A^{\left(S\right)},S^{\prime}\mapsto A^{\left(S^{\prime}\right)}$是典范嵌入.
\end{corollary}

\begin{proposition}{}{}
	设$T$是$S$的子原群,则$\left\{\sum\limits_{s\in T}\alpha_{s}e_{s}\in A^{\left(S\right)}\middle|\left(\alpha_{s}\right)_{s\in T}\text{有限支撑}\right\}$是$A^{\left(S\right)}$的子代数,并且典范同构于$A^{\left(T\right)}$.
\end{proposition}

\begin{example}{幺半群作用到模的转换}{}
	设$V$是$A$-模,$S$是幺半群且左作用于$M$,于是:
	\begin{itemize}
		\item 上述左作用给出$f\in\Hom_{\mathsf{Mon}}\left(S,\End_{A}\left(V\right)\right)$,它满足$f\left(1_{S}\right)=\id_{V}$.
		\item $f$诱导$\bar{f}\in\Hom_{A-\mathsf{Alg}}\left(A^{\left(S\right)},\End_{A}\left(V\right)\right)$.
	\end{itemize}
	因此,$V$是典范的左$A^{\left(S\right)}$-模.
\end{example}

\begin{example}{$\Omega$-群到模的转换}{}
	设$M$是带算子的交换群($M$的运算采用加法记号,算子作用采用乘法记号),$\Omega$是这族算子的定义域的无交并,将每个算子的定义域等同于$\Omega$的一个子集.定义$M$在$\Mo\left(\Omega\right)$上的左作用为$\left(s,x\right)\mapsto s\cdot x$,则$M$是典范的$\Z^{\left(\Mo\left(\Omega\right)\right)}$-模.因此,关于带算子的交换群的性质和与之相应的模的概念是一致的.
\end{example}












































